\documentclass[11pt,a4paper]{article}
\usepackage[T1]{fontenc}
\usepackage[utf8]{inputenc}
\usepackage[provide=*,lithuanian]{babel}
\usepackage{lmodern,amsmath,amssymb,enumitem}
\usepackage[margin=1.4cm]{geometry}
\setlength{\parindent}{0pt}
\setlist[enumerate]{label=\textbf{\arabic*.},leftmargin=0.65cm,itemsep=0.35cm,topsep=0.2cm,parsep=0pt}
\begin{document}
\begin{center}
{\Large\bfseries Lygtys su keitiniais}

\medskip
{\large Lygis: Beast}

\smallskip
30 mišrių uždavinių
\end{center}
\textbf{Užduotis.} Pertvarkykite lygtį, įveskite tinkamą keitinį \(t\),
išspręskite gautą pilnąją kvadratinę lygtį ir raskite pradinės lygties
realiuosius sprendinius. Atsižvelkite į pradinių reiškinių apibrėžimo sritį
ir galimas keitinio reikšmes.
\medskip
\begin{enumerate}
\item \(\displaystyle x^2(x-3)^2-2x^2+6x-8=0\)
\item \(\displaystyle (x^2-x+2)^2=2x(x-1)+12\)
\item \(\displaystyle 3^{x+1}+3^{1-x}=10\)
\item \(\displaystyle x^2(x+2)^2-5x^2-10x+6=0\)
\item \(\displaystyle (x^2+3x+1)^2=(x+2)(x+1)+19\)
\item \(\displaystyle 2^{x+1}+2^{1-x}=5\)
\item \(\displaystyle \frac{2x+6}{x-1}+\frac{6x-6}{2x+6}=7\)
\item \(\displaystyle (x^2+2x-1)^2=3(x+1)^2+4\)
\item \(\displaystyle 4x^2(x+1)^2-6x^2-6x+2=0\)
\item \(\displaystyle 5^{x+1}+5^{2-x}=30\)
\item \(\displaystyle \frac{3x-6}{x+1}+\frac{2x+2}{x-2}=7\)
\item \(\displaystyle (x^2+2)^2-4x(x^2+2)+4x^2-7x^2+14x-4=0\)
\item \(\displaystyle (2x^2+3x+1)^2=(2x+1)(x+1)+6\)
\item \(\displaystyle 2^{x+2}+2^{3-x}=18\)
\item \(\displaystyle \frac{x^2+3x+2}{x^2-1}-\frac{6x-6}{x+2}=1\)
\item \(\displaystyle (x^2+x)^2+2(x^2+x)(x-1)+(x-1)^2-5x^2-10x+11=0\)
\item \(\displaystyle (x-1)^2(x+2)^2-2x^2-2x-4=0\)
\item \(\displaystyle 3^{x+1}-3^{2-x}=26\)
\item \(\displaystyle (x^2+2x)^2+2(x^2+2x)(x^2-x+1)+(x^2-x+1)^2-10x^2-5x-1=0\)
\item \(\displaystyle \frac{2x^2+2}{x^2-1}-\frac{3x^2-3}{x^2+1}=1\)
\item \(\displaystyle (x+1)^2(x-2)^2-3x^2+3x+2=0\)
\item \(\displaystyle (2x^2+x+1)^2-2(2x^2+x+1)(x^2+3x)+(x^2+3x)^2-5x^2+10x-1=0\)
\item \(\displaystyle (x^2+x-3)^2+x(x+1)=9\)
\item \(\displaystyle \frac{x^2-x-2}{x^2-4}-\frac{2x+4}{x+1}=1\)
\item \(\displaystyle (x^2-4x+5)^2=2(x-1)(x-3)+19\)
\item \(\displaystyle 2x^2(2x-1)^2-6x^2+3x+1=0\)
\item \(\displaystyle (x^2-1)^2+(x+1)^2+2(x^2-1)(x+1)-8x^2-8x+12=0\)
\item \(\displaystyle \left(\frac12\right)^{x-1}+\left(\frac12\right)^{-x}=3\)
\item \(\displaystyle \frac{4x^2+4x+4}{2x^2+2x-2}+\frac{3x^2+3x-3}{x^2+x+1}=7\)
\item \(\displaystyle (x^2+x+2)^2+(x^2-2x+1)^2-2(x^2+x+2)(x^2-2x+1)-3x-7=0\)
\end{enumerate}
\clearpage
\section*{Atsakymai}
\begingroup\small
\begin{enumerate}
\item \(\displaystyle x_{1}=-1,\quad x_{2}=1,\quad x_{3}=2,\quad x_{4}=4\)
\item \(\displaystyle x_{1}=-1,\quad x_{2}=2\)
\item \(\displaystyle x_{1}=-1,\quad x_{2}=1\)
\item \(\displaystyle x_{1}=-3,\quad x_{2}=-1-\sqrt{3},\quad x_{3}=-1+\sqrt{3},\quad x_{4}=1\)
\item \(\displaystyle x_{1}=-4,\quad x_{2}=1\)
\item \(\displaystyle x_{1}=-1,\quad x_{2}=1\)
\item \(\displaystyle x_{1}=-7,\quad x_{2}=3\)
\item \(\displaystyle x_{1}=-1-\sqrt{7},\quad x_{2}=-1,\quad x_{3}=-1+\sqrt{7}\)
\item \(\displaystyle x_{1}=\frac{-1-\sqrt{5}}{2},\quad x_{2}=\frac{-1-\sqrt{3}}{2},\quad x_{3}=\frac{-1+\sqrt{3}}{2},\quad x_{4}=\frac{-1+\sqrt{5}}{2}\)
\item \(\displaystyle x_{1}=0,\quad x_{2}=1\)
\item \(\displaystyle x_{1}=-4,\quad x_{2}=\frac{7}{2}\)
\item \(\displaystyle x_{1}=-1,\quad x_{2}=0,\quad x_{3}=2,\quad x_{4}=3\)
\item \(\displaystyle x_{1}=-2,\quad x_{2}=\frac{1}{2}\)
\item \(\displaystyle x_{1}=-1,\quad x_{2}=2\)
\item \(\displaystyle x_{1}=0,\quad x_{2}=\frac{5}{2}\)
\item \(\displaystyle x_{1}=-1-\sqrt{5},\quad x_{2}=-3,\quad x_{3}=1,\quad x_{4}=-1+\sqrt{5}\)
\item \(\displaystyle x_{1}=-3,\quad x_{2}=-1,\quad x_{3}=0,\quad x_{4}=2\)
\item \(\displaystyle x_{1}=2\)
\item \(\displaystyle x_{1}=-\frac{3}{2},\quad x_{2}=-\frac{1}{2},\quad x_{3}=0,\quad x_{4}=1\)
\item \(\displaystyle x_{1}=-\sqrt{5},\quad x_{2}=0,\quad x_{3}=\sqrt{5}\)
\item \(\displaystyle x_{1}=-2,\quad x_{2}=\frac{1-\sqrt{5}}{2},\quad x_{3}=\frac{1+\sqrt{5}}{2},\quad x_{4}=3\)
\item \(\displaystyle x_{1}=-1,\quad x_{2}=0,\quad x_{3}=2,\quad x_{4}=3\)
\item \(\displaystyle x_{1}=\frac{-1-\sqrt{21}}{2},\quad x_{2}=-1,\quad x_{3}=0,\quad x_{4}=\frac{-1+\sqrt{21}}{2}\)
\item \(\displaystyle x_{1}=-3,\quad x_{2}=-\frac{3}{2}\)
\item \(\displaystyle x_{1}=0,\quad x_{2}=4\)
\item \(\displaystyle x_{1}=-\frac{1}{2},\quad x_{2}=\frac{1-\sqrt{5}}{4},\quad x_{3}=\frac{1+\sqrt{5}}{4},\quad x_{4}=1\)
\item \(\displaystyle x_{1}=-3,\quad x_{2}=-2,\quad x_{3}=1,\quad x_{4}=2\)
\item \(\displaystyle x_{1}=0,\quad x_{2}=1\)
\item \(\displaystyle x_{1}=-2,\quad x_{2}=1\)
\item \(\displaystyle x_{1}=-1,\quad x_{2}=\frac{2}{3}\)
\end{enumerate}
\endgroup
\clearpage
\section*{Keitinių užuominos}
\begin{enumerate}
\item \(\displaystyle t=x(x-3)\)
\item \(\displaystyle t=x^2-x+2\)
\item \(\displaystyle t=3^x\)
\item \(\displaystyle t=x(x+2)\)
\item \(\displaystyle t=x^2+3x+1\)
\item \(\displaystyle t=2^x\)
\item \(\displaystyle t=\frac{x+3}{x-1}\)
\item \(\displaystyle t=x^2+2x-1\)
\item \(\displaystyle t=2x(x+1)\)
\item \(\displaystyle t=5^x\)
\item \(\displaystyle t=\frac{x-2}{x+1}\)
\item \(\displaystyle t=x^2-2x+2\)
\item \(\displaystyle t=2x^2+3x+1\)
\item \(\displaystyle t=2^x\)
\item \(\displaystyle t=\frac{x+2}{x-1}\)
\item \(\displaystyle t=x^2+2x-1\)
\item \(\displaystyle t=(x-1)(x+2)\)
\item \(\displaystyle t=3^x\)
\item \(\displaystyle t=2x^2+x+1\)
\item \(\displaystyle t=\frac{x^2+1}{x^2-1}\)
\item \(\displaystyle t=(x+1)(x-2)\)
\item \(\displaystyle t=x^2-2x+1\)
\item \(\displaystyle t=x^2+x-3\)
\item \(\displaystyle t=\frac{x+1}{x+2}\)
\item \(\displaystyle t=x^2-4x+5\)
\item \(\displaystyle t=x(2x-1)\)
\item \(\displaystyle t=x^2+x\)
\item \(\displaystyle t=\left(\frac12\right)^x\)
\item \(\displaystyle t=\frac{x^2+x+1}{x^2+x-1}\)
\item \(\displaystyle t=3x+1\)
\end{enumerate}
\end{document}
